\documentclass[11pt]{article}

\usepackage[T1]{fontenc}
\usepackage[utf8]{inputenc}
\usepackage{lmodern}
\usepackage{microtype}
\usepackage[margin=1in]{geometry}
\usepackage{amsmath}
\usepackage{amssymb}
\usepackage{amsthm}
\usepackage{booktabs}
\usepackage{enumitem}
\usepackage{xcolor}
\usepackage{url}
\usepackage[
  colorlinks=true,
  linkcolor=blue!55!black,
  citecolor=blue!55!black,
  urlcolor=blue!55!black,
  pdftitle={A Level-15 Certificate for the Prime-29 Gate in the Fifteen-Runner Lonely Runner Problem},
  pdfauthor={Ruturaj R Raval},
  pdfsubject={Computer-assisted finite checking for the Lonely Runner Conjecture},
  pdfkeywords={lonely runner conjecture, Diophantine approximation, SAT, DRAT, exact computation}
]{hyperref}

\newtheorem{theorem}{Theorem}
\newtheorem{lemma}[theorem]{Lemma}
\newtheorem{proposition}[theorem]{Proposition}
\newtheorem{corollary}[theorem]{Corollary}
\newtheorem{remark}[theorem]{Remark}

\newcommand{\distm}[2]{\operatorname{dist}_{#1}(#2)}
\newcommand{\Z}{\mathbb{Z}}
\newcommand{\Jset}{J(14,29)}
\setlength{\emergencystretch}{2em}

\title{A Level-15 Certificate for the Prime-29 Gate\\
in the Fifteen-Runner Lonely Runner Problem}
\author{%
  Ruturaj R Raval\\
  \small Independent Researcher\\
  \small ORCID:
  \href{https://orcid.org/0000-0003-4930-8981}
       {\texttt{0000-0003-4930-8981}}
}
\date{September 6, 2026}

\begin{document}

\maketitle

\begin{abstract}
The finite-checking approach to the Lonely Runner Conjecture associates a
set \(J(k,p)\) of eventually improper tuples with each prime \(p\).  This
report proves
\[
  J(14,29)=\varnothing,
\]
closing the prime-29 gate for fifteen total runners.  At level one modulo
29, every folded speed class is bad at exactly one folded time class.  The
unique improper orbit is therefore represented by
\((1,2,\ldots,14)\).

Every improper level-15 lift of this orbit is equivalent to a lift in one of
two exhaustive symmetry branches.  In the branch containing a coordinate
coprime to 15, a
Chinese-remainder reduction forces a zero residue modulo 15 in one of
coordinates 2 through 14.  Separately written C++ and Rust exact solvers
reject all thirteen resulting cases.  In the branch with no coordinate
coprime to 15, a 210-variable, 1842-clause CNF is unsatisfiable.  A
368542-byte DRAT proof is accepted by both DRAT-trim and \texttt{rate}.

The release certificate binds the source, formula, proof, logs, and original
runs by SHA-256 and supports full recompilation and replay.  The closed gate
contributes \(\log 29\) to the finite-checking prime mass.  It does not prove
the complete fifteen-runner conjecture.
\end{abstract}

\section{Introduction}

The Lonely Runner Conjecture asks whether, for \(n\) runners with distinct
constant speeds on a unit circle, every runner is at some time at distance
at least \(1/n\) from every other runner.  After subtracting the speed of one
runner, the problem is expressed using \(k=n-1\) nonzero integer speeds.  We
write \(LRC(k)\) for this \(k\)-moving-runner form.

Sungkawichai and Trakulthongchai developed a finite-checking method based on
improper modular families, exact lifting, backward projection, and prime
gates \cite{sungkawichai-trakulthongchai}.  Their work establishes the
conjecture through thirteen total runners.  Allikvere subsequently reported
a computer-assisted proof for fourteen total runners
\cite{allikvere-fourteen}.  The next full case is therefore \(LRC(14)\),
corresponding to fifteen total runners.

For each prime \(p\), the finite-checking method defines an
eventual-improper set \(J(k,p)\).  Proving \(J(k,p)=\varnothing\) closes the
prime-\(p\) gate.  A sufficient collection of closed gates proves \(LRC(k)\)
when the sum of their logarithms exceeds an explicit threshold.  For
\(k=14\), the current threshold is
\[
  \log B_{14}=810.0739811140556.
\]

This report closes one gate:
\[
  J(14,29)=\varnothing.
\]
The contribution is deliberately scoped.  Prime 29 contributes only
\[
  \log 29=3.367295829986474,
\]
about \(0.41568\) percent of the threshold.  The result is a certified input
to a future full proof, not that proof itself.

\section{Finite gate definitions}

Fix \(k\geq 2\), put \(q=k+1\), and let \(p\) be prime.  At level
\(\ell\geq 1\), the modulus is \(m=\ell p\).  For \(x\in\Z\), define
\[
  \distm{m}{x}=\min(x\bmod m,-x\bmod m),
\]
where both residues are represented in \(\{0,\ldots,m-1\}\).

A \(k\)-tuple \(v=(v_1,\ldots,v_k)\), with no coordinate divisible by
\(p\), has a lonely time witness at level \(\ell\) if some time class
\(t\bmod m\) satisfies
\[
  q\,\distm{m}{t v_i}\geq m
  \qquad\text{for every }i.
\]
There is also a gcd properness condition: the tuple is proper if, for some
coordinate \(i\),
\[
  \gcd(\ell,v_1,\ldots,v_{i-1},v_{i+1},\ldots,v_k)>1.
\]
A tuple is improper when neither condition holds.

A level-one tuple is eventually proper if some finite level contains no
improper lift of that tuple.  The level-one tuples that are not eventually
proper form \(J(k,p)\).  Coordinate permutations, independent sign changes,
and multiplication by units preserve eventual properness.

For the remainder of the report,
\[
  k=14,\qquad q=15,\qquad p=29.
\]
Folding by sign leaves fourteen nonzero speed classes and fourteen nonzero
time classes modulo 29.

\section{The unique level-one orbit}

For a folded speed class \(v\in\{1,\ldots,14\}\), define its bad-time set
\[
  B_v=\{t\in\{1,\ldots,14\}:
        15\,\distm{29}{t v}<29\}.
\]

\begin{lemma}
\label{lem:singleton-bad}
Every \(B_v\) contains exactly one folded time class, and the sets
\(B_1,\ldots,B_{14}\) partition the folded nonzero time classes.
\end{lemma}

\begin{proof}
For nonzero \(t,v\bmod 29\), the product \(tv\) is nonzero.  The strict
inequality
\[
  15\,\distm{29}{tv}<29
\]
therefore holds exactly when \(\distm{29}{tv}=1\), or equivalently
\(tv\equiv\pm1\pmod{29}\).  For fixed \(v\), this determines one folded
time class.  Multiplication by \(v^{-1}\) permutes the nonzero classes, so
the fourteen singleton sets partition the folded time classes.
\end{proof}

\begin{proposition}
\label{prop:level-one}
Up to permutation and independent signs, the unique level-one improper
14-tuple is
\[
  (1,2,\ldots,14).
\]
\end{proposition}

\begin{proof}
A level-one tuple is improper only if the union of the bad-time sets of its
coordinates covers all fourteen folded time classes.  By Lemma
\ref{lem:singleton-bad}, each coordinate covers one class.  Fourteen
coordinates cover all fourteen classes only when each folded speed class
occurs exactly once.
\end{proof}

The release also contains an exact level-one generator.  It enumerates every
covering support, expands multiplicities to total size 14, and quotients by
permutation, signs, and units.  Its output is the same single orbit, with
canonical row hash
\[
\texttt{a8d0f6ffa5961efb0405e27450d85cef2b471d9dcb4d2afc73b5f71c7127ed08}.
\]

\section{Level-15 lifts and symmetry}

Label the coordinates of the representative from Proposition
\ref{prop:level-one}.  Every level-15 lift can be written
\[
  v_i=i+29a_i,
  \qquad 1\leq i\leq14,\quad 0\leq a_i<15.
\]
Define
\[
  c_i\equiv v_i\equiv i-a_i\pmod{15}.
\]
For each coordinate, \(a_i\mapsto c_i\) is a bijection of
\(\Z/15\Z\).

\begin{lemma}
\label{lem:symmetry-split}
Every improper level-15 lift is equivalent to a lift in one of the following
branches.
\begin{enumerate}[label=(\alph*),leftmargin=2.2em]
  \item The coprime branch, normalized by \(v_1=1\), hence \(c_1=1\).
  \item The noncoprime branch, in which every coordinate is noncoprime to
        15 and \(v_3=3\).
\end{enumerate}
\end{lemma}

\begin{proof}
Suppose first that some coordinate is coprime to 15.  It is also nonzero
modulo 29 and is therefore a unit modulo \(435=15\cdot29\).
Multiplication by its inverse, followed by a sign change and coordinate
permutation, normalizes that coordinate to 1.  Since the parent contains
every nonzero folded class modulo 29, the normalized parent remains in the
same orbit.  We label the normalized coordinate as coordinate 1.

Now suppose no coordinate is coprime to 15.  An improper tuple must fail the
gcd properness condition after every coordinate deletion.  In particular,
at least two coordinates are not divisible by 3 and at least two are not
divisible by 5.  A coordinate not divisible by 5 cannot be coprime to 15 in
this branch, so its gcd with 15 is exactly 3.  Such a coordinate can be
normalized to 3 by a unit modulo 435, a sign change, and a permutation.
Unit multiplication preserves every coordinate's gcd with 15.
\end{proof}

\section{CRT reduction of the coprime branch}

Let \(1\leq t\leq217\) represent a folded nonzero time class modulo 435.
Put
\[
  x\equiv t\pmod{29},\qquad
  y\equiv t\pmod{15},\qquad
  r\equiv xi\pmod{29},
\]
with \(r\in\{0,\ldots,28\}\).

\begin{lemma}
\label{lem:crt}
Coordinate \(i\) is bad at time \(t\) exactly under the following
conditions:
\[
  \begin{cases}
  yc_i\equiv0\pmod{15}, & r=0,\\
  yc_i\equiv r\text{ or }r+1\pmod{15}, & r\neq0.
  \end{cases}
\]
\end{lemma}

\begin{proof}
Badness means
\[
  15\,\distm{435}{t v_i}<435,
\]
so the signed representative of \(tv_i\bmod435\) has absolute value at most
28.  Its residue modulo 29 is \(r\).

If \(r=0\), the only representative in the allowed interval is 0.  The
residue modulo 15 must therefore be 0.

If \(r\neq0\), the two representatives in the allowed interval with residue
\(r\bmod29\) are \(r\) and \(r-29\).  Their residues modulo 15 are \(r\)
and \(r+1\), because \(-29\equiv1\pmod{15}\).  The residue of \(tv_i\)
modulo 15 is \(yc_i\), proving the claim.
\end{proof}

Three consequences make the search finite and small enough for exhaustive
replay.

\begin{enumerate}[leftmargin=2.2em]
  \item At \(t=29\), Lemma \ref{lem:crt} says that a coordinate is bad
        exactly when \(c_i=0\).  Every time cover contains a zero residue.
  \item Every folded multiple of 15 is covered automatically.  At each such
        time, one parent coordinate is bad for every value of its residue
        \(c_i\).
  \item A coordinate with \(c_i=0\) covers every folded multiple of 29.
\end{enumerate}

There are 217 folded time classes, fourteen multiples of 15, and seven
multiples of 29.  No class in this range is a multiple of both, leaving
\[
  217-14-7=196
\]
nontrivial time classes.  The normalization \(c_1=1\) puts the mandatory
zero in one of coordinates \(2,\ldots,14\).  Thus the thirteen cases
\[
  c_j=0,\qquad j=2,\ldots,14,
\]
are exhaustive.  They overlap when a tuple has more than one zero, which is
harmless for an unsatisfiability proof.

\section{Exact coprime searches}

Two separately written programs solve all thirteen coordinate-zero cases.

The C++ implementation evaluates badness directly modulo 435 and branches
on coordinate groups.  The Rust implementation uses Lemma \ref{lem:crt}
and branches on individual clause literals.  Before the full replay launches
the Rust searches, it invokes the binary's explicit self-test, which compares
the CRT predicate exhaustively with direct modular arithmetic.

Both programs omit the gcd properness clauses in this branch.  They prove
the stronger statement that no normalized lift even covers every folded
time class.

\begin{table}[htbp]
\centering
\caption{Aggregate exact-search records for the thirteen coprime cases.}
\label{tab:coprime-search}
\begin{tabular}{lrrr}
\toprule
Implementation & Search nodes & Propagations & Fallback branches\\
\midrule
C++ direct arithmetic & 58,222,690 & 66,402,789 & 0\\
Rust CRT              & 1,319,160,790 & 1,654,750,357 & 0\\
\bottomrule
\end{tabular}
\end{table}

Every case exits with the exact status \texttt{UNSAT}.  The node counts need
not agree because the branching units and arithmetic representations differ.
The packaged full-replay command recompiles both source snapshots and reruns
all twenty-six computations.

\begin{proposition}
\label{prop:coprime}
No improper level-15 lift exists in the coprime branch.
\end{proposition}

\begin{proof}
Every lift in the branch that covers all folded times has \(c_j=0\) for at
least one \(j\in\{2,\ldots,14\}\).  Both exact implementations exhaust every
assignment in each of these thirteen cases and return \texttt{UNSAT}.
Therefore no assignment covers all folded times.
\end{proof}

\section{The noncoprime DRAT certificate}

The noncoprime branch is encoded as a Boolean formula with one variable for
each of the fifteen choices of each of fourteen labeled coordinates, giving
210 variables.  The formula contains:

\begin{itemize}[leftmargin=2em]
  \item exactly-one constraints for each coordinate;
  \item one bad-coordinate clause for every folded time class;
  \item the exact simultaneous failure conditions for the factor-3 and
        factor-5 gcd properness clauses;
  \item unit clauses excluding choices coprime to 15; and
  \item the symmetry anchor \(v_3=3\).
\end{itemize}

The resulting DIMACS instance has 1842 clauses.  Kissat 4.0.4 returns
\texttt{UNSATISFIABLE} and produces a DRAT proof.  The principal artifact
hashes are:

\begin{center}
\begin{tabular}{lll}
\toprule
Artifact & Size & SHA-256\\
\midrule
CNF & 52,466 bytes &
\texttt{9869360a6aab0a6...d5b9b5fc5}\\
DRAT proof & 368,542 bytes &
\texttt{1554d825fa23b8b...933c9bacd}\\
Summary & 3,737 bytes &
\texttt{937d711941b68c2...344e4e16d}\\
\bottomrule
\end{tabular}
\end{center}

The complete hashes appear in the machine-readable manifest.  The checker
source revisions are:
\begin{center}
\small
\begin{tabular}{ll}
DRAT-trim &
\texttt{2e3b2dc0ecf938addbd779d42877b6ed69d9a985}\\
\texttt{rate} &
\texttt{93bf524435455ddc0b370669c17b1eb8741a026c}\\
\end{tabular}
\end{center}
Both checkers accept the proof with the exact line
\texttt{s VERIFIED}.

\begin{proposition}
\label{prop:noncoprime}
No improper level-15 lift exists in the noncoprime branch.
\end{proposition}

\begin{proof}
The CNF is an exact encoding of the normalized branch and of simultaneous
failure of both properness mechanisms.  Its DRAT refutation is accepted by
two independently implemented checkers.  Therefore the formula is
unsatisfiable.
\end{proof}

\section{Main theorem}

\begin{theorem}
\label{thm:main}
The prime-29 eventual-improper set for fourteen moving runners is empty:
\[
  J(14,29)=\varnothing.
\]
\end{theorem}

\begin{proof}
By Proposition \ref{prop:level-one}, every level-one improper tuple belongs
to the orbit represented by \((1,2,\ldots,14)\).  Lemma
\ref{lem:symmetry-split} puts every improper level-15 lift into the coprime
or noncoprime branch.  The first is excluded by Proposition
\ref{prop:coprime}; the second is excluded by Proposition
\ref{prop:noncoprime}.  Thus the unique level-one orbit has no improper
level-15 lift and is eventually proper.
\end{proof}

\section{Certificate and trust boundary}

The release package separates three checks:

\begin{enumerate}[leftmargin=2.2em]
  \item exact regeneration of the unique level-one orbit;
  \item dual-language exhaustive rejection of the coprime cases; and
  \item proof-producing SAT rejection of the noncoprime case.
\end{enumerate}

Every included artifact is listed by path, byte count, and SHA-256 in a
top-level manifest.  A bundled Git object archive binds the exact source
history to the declared commit.  The fast verifier checks out that commit,
compares every packaged solver and generator source with its Git object, and
runs the original reconstruction verifier from the checked-out source.  The
three original runs share source commit
\[
\texttt{f2f933bf16ab8eb5b081ddca59b18bf6b1b92e9d}.
\]
Each production runner verifies a clean tracked worktree, an unchanged
40-character commit, and every source, executable, runner, Makefile, or
proof-tool hash declared in its own provenance record before accepting its
output.

The standalone verifier has three modes.

\begin{description}[leftmargin=2.8em,style=nextline]
  \item[Fast integrity and reconstruction]
    Checks the complete artifact manifest, exact claim fields, source
    bundle, all transcripts, level-one generation, CRT identities, and a
    fresh reconstruction of the noncoprime CNF.  It does not execute the
    long searches or proof checkers.
  \item[DRAT replay]
    Runs DRAT-trim and \texttt{rate} on the packaged formula and proof.
  \item[Full coprime replay]
    Compiles the C++ and Rust sources from inside the certificate and reruns
    all thirteen cases with each implementation.
\end{description}

Adversarial tests alter the theorem field, substitute the source commit,
corrupt the Git bundle, add an unlisted nested \texttt{certificate.json},
change packaged source while refreshing the outer artifact hash, and corrupt
a structured summary.  Each mutation is rejected without an uncaught
traceback.

The noncoprime branch has a compact standard proof object.  The coprime
branch does not.  Its strongest assurance comes from source review,
different arithmetic formulations, fail-closed transcripts, and complete
recompilation and replay.

\section{Significance and remaining work}

Theorem \ref{thm:main} contributes a certified prime gate to the
finite-checking program.  It also supplies a factor-15 CRT reduction that can
be tested on other primes and a source-bound release format suitable for long
exact searches.

The quantitative gap remains large.  Prime 29 contributes about
\(0.41568\) percent of the required logarithmic mass.  A complete proof of
\(LRC(14)\) requires many additional gates or a structural theorem that
closes a substantial family at once.

The immediate research directions are:

\begin{itemize}[leftmargin=2em]
  \item rank additional primes by level-one orbit count and forced-residue
        structure;
  \item extend the CRT analysis to those primes;
  \item use proof-producing SAT whenever a branch admits a compact encoding;
  \item retain separately written exhaustive solvers when a compact proof
        object is unavailable; and
  \item track cumulative certified prime mass against \(\log B_{14}\).
\end{itemize}

No claim is made that the dated public-source search found unpublished or
unindexed computations.  No external peer review is claimed.

\begingroup
\raggedright
\begin{thebibliography}{9}

\bibitem{sungkawichai-trakulthongchai}
T.~Sungkawichai and T.~Trakulthongchai,
\emph{Eleven, twelve, and thirteen lonely runners},
arXiv:2604.23906v1, 2026.
\url{https://arxiv.org/abs/2604.23906}

\bibitem{allikvere-fourteen}
J.~Allikvere,
\emph{Fourteen lonely runners},
arXiv:2609.02604v1, 2026.
\url{https://arxiv.org/abs/2609.02604}

\bibitem{allikvere-archive}
J.~Allikvere,
\emph{Verification archive for Fourteen lonely runners},
Zenodo, 2026.
\url{https://doi.org/10.5281/zenodo.22066772}

\end{thebibliography}
\endgroup

\end{document}
